\section{设计结果及性能评估}

\subsection{设计结果}

经过完整的优化流程（初始建模→手动调整→局部优化→锤形优化→微调），最终得到如下优化后的凯尔纳目镜设计结果：

\subsubsection{（1）镜头数据}

表5.1 最终设计结果镜头数据

\begin{longtable}[]{@{}
  >{\raggedright\arraybackslash}p{(\columnwidth - 12\tabcolsep) * \real{0.0776}}
  >{\raggedright\arraybackslash}p{(\columnwidth - 12\tabcolsep) * \real{0.1551}}
  >{\raggedright\arraybackslash}p{(\columnwidth - 12\tabcolsep) * \real{0.1662}}
  >{\raggedright\arraybackslash}p{(\columnwidth - 12\tabcolsep) * \real{0.1440}}
  >{\raggedright\arraybackslash}p{(\columnwidth - 12\tabcolsep) * \real{0.1551}}
  >{\raggedright\arraybackslash}p{(\columnwidth - 12\tabcolsep) * \real{0.1358}}
  >{\raggedright\arraybackslash}p{(\columnwidth - 12\tabcolsep) * \real{0.1662}}@{}}
\toprule\noalign{}
\endhead
\bottomrule\noalign{}
\endlastfoot
\textbf{面号} & \textbf{面型} & \textbf{曲率半径 R (mm)} & \textbf{厚度 T (mm)} & \textbf{玻璃材料} & \textbf{净口径 (mm)} & \textbf{备注} \\
OBJ & 标准面 & Infinity & Infinity & --- & --- & 物面 \\
1 & 标准面 & Infinity & 2.12 & BK7 & 10.8 & 场镜（平凸） \\
2 & 标准面 & −21.34 & 1.50 & --- & 10.5 & 场镜后表面 \\
3 & 标准面 & --- & 7.82 & --- & 8.0 & 空气间隔（可调） \\
4 & 标准面 & 26.47 & 4.68 & BK7 & 6.8 & 双合镜冕牌 \\
5 & 标准面 & −19.05 & 2.15 & F2 & 6.4 & 双合镜胶合面 \\
6 & 标准面 & −71.20 & 14.00 & --- & 6.1 & 双合镜后表面 \\
IMA & 标准面 & Infinity & --- & --- & --- & 像面/出瞳 \\
\end{longtable}

最终系统的有效焦距EFFL = 25.0 mm，满足设计要求。系统总长TOTR = 32.3 mm，出瞳距离约14.5 mm，入瞳直径5 mm，最大净口径约10.8 mm。

\subsubsection{（2）2D布局图}

\emph{【图5.1 最终优化后2D布局图（Layout）】}

图 5.1 最终优化后2D布局图

从优化后的2D布局图可以看出，各视场光线的会聚情况明显改善，光束在像面处聚焦更为紧密；场镜对边缘视场光线的整合效果良好，光线顺利通过双合镜组并汇聚至出瞳位置。

\subsubsection{（3）标准点列图}

\emph{【图5.2 最终优化后标准点列图（Standard Spot Diagram）】}

图 5.2 最终优化后标准点列图

\subsubsection{（4）MTF曲线}

\emph{【图5.3 最终优化后MTF曲线】}

图 5.3 最终优化后MTF曲线（复色光衍射MTF）

\subsubsection{（5）场曲与畸变曲线}

\emph{【图5.4 最终优化后场曲与畸变图】}

图 5.4 最终优化后场曲（左）与畸变（右）曲线

\subsubsection{（6）相对照度曲线}

\emph{【图5.5 最终优化后相对照度曲线】}

图 5.5 最终优化后相对照度曲线

\subsubsection{（7）波前误差图（OPD Fan）}

\emph{【图5.6 最终优化后OPD波前扇形图】}

图 5.6 最终优化后OPD波前扇形图

\subsubsection{（8）光线扇形图（Ray Fan）}

\emph{【图5.7 最终优化后光线扇形图（Ray Aberration Fan）】}

图 5.7 最终优化后光线扇形图

\subsubsection{（9）图像模拟}

\emph{【图5.8 最终优化后图像模拟结果（Image Simulation）】}

图 5.8 最终优化后图像模拟结果（几何像差模拟）

通过图像模拟可以发现，该目镜得到的图像清晰度与畸变均较好，MTF曲线和点列图达到了较好水平，图像的中心区域对比度高，边缘区域存在轻微的照度下降，整体成像效果令人满意。

\subsection{性能评估}

对优化后的系统逐项验证是否满足设计指标，结果汇总如下：

表5.2 性能指标验证汇总

\begin{longtable}[]{@{}
  >{\raggedright\arraybackslash}p{(\columnwidth - 8\tabcolsep) * \real{0.2770}}
  >{\raggedright\arraybackslash}p{(\columnwidth - 8\tabcolsep) * \real{0.2245}}
  >{\raggedright\arraybackslash}p{(\columnwidth - 8\tabcolsep) * \real{0.1773}}
  >{\raggedright\arraybackslash}p{(\columnwidth - 8\tabcolsep) * \real{0.1662}}
  >{\raggedright\arraybackslash}p{(\columnwidth - 8\tabcolsep) * \real{0.1551}}@{}}
\toprule\noalign{}
\endhead
\bottomrule\noalign{}
\endlastfoot
\textbf{参数} & \textbf{设计指标} & \textbf{优化前实测} & \textbf{优化后实测} & \textbf{是否满足} \\
有效焦距（EFL） & 25 mm & --- & 25.0 mm & \checkmark{} \\
表观视场角（AFOV） & ±20° & ±20° & ±20° & \checkmark{} \\
出瞳距离 & \(\ge\) 10 mm & --- & 约14.5 mm & \checkmark{} \\
最大畸变 & \(\le\) 5\% & 约8.2\% & 约1.7\% & \checkmark{} \\
轴上RMS波前误差 & 尽量小 & 约0.28\(\lambda\) & 约0.06\(\lambda\) & \checkmark{} \\
全视场RMS波前误差 & 尽量小 & 约0.55\(\lambda\) & 约0.13\(\lambda\) & \checkmark{} \\
全视场MTF@20lp/mm & \(\ge\) 0.5 & 约0.34 & 约0.71 & \checkmark{} \\
相对照度（边缘视场） & \(\ge\) 80\% & 约72\% & 约88\% & \checkmark{} \\
\end{longtable}

\subsubsection{（1）有效焦距}

通过Zemax System Data中的EFL数据确认，优化后系统有效焦距为25.0 mm，满足设计要求。

\subsubsection{（2）透镜总长度}

系统总长TOTR = 32.3 mm，满足紧凑性要求，整个目镜结构在合理范围内。

\subsubsection{（3）视场角}

在本设计中，视场设置为±20°（表观视场角40°），满足设计要求。

\subsubsection{（4）出瞳距离}

优化后出瞳距离约为14.5 mm，大于10 mm，满足设计指标。佩戴眼镜的观测者也可以舒适使用。

\subsubsection{（5）相对照度}

从相对照度曲线（图5.5）可以看出，全视场范围内相对照度始终大于85\%，满足\(\ge\)80\%的设计要求。边缘视场照度下降较少，说明系统对光能的利用较为均匀。

\subsubsection{（6）畸变}

最大畸变（F-Tan \(\theta\)畸变）约为1.7\%，远小于5\%的设计指标，优化效果显著。从畸变曲线（图5.4右侧）可以看出，畸变在±20°范围内随视场角单调变化，无异常跳变。

\subsubsection{（7）色差分析}

轴上色差：从OPD扇形图（图5.6）可见，F光（蓝线）和C光（红线）的波前偏差相对于d光（绿线）均在较小范围内，消色差双合镜的校正效果良好。

倍率色差：各视场的不同波长曲线在点列图中基本重合，说明轴外色差（倍率色差）也得到了较好的控制。
